% Direct-authority diplomatic French for printed p.150.
\SourcePage{150}

\LayerTarget{tohoku:prop:2.5.2:proof:p2}
\noindent dans le diagramme analogue relatif à $C'(A)$, ce qui nous montre
que l'isomorphisme obtenu est indépendant du choix de $C(A)$. D'autre part,
se bornant désormais à des résolutions injectives, on voit immédiatement que
les isomorphismes obtenus sont fonctoriels et permutent aux homomorphismes
cobords, de plus ils se réduisent en dimension 0 à l'isomorphisme
d'augmentation (ou plutôt son inverse), donc ces isomorphismes sont bien ceux
de la proposition 1.

\LayerTarget{tohoku:example:2.5:identity-resolution}
\textsc{Exemples.} a) Considérons le foncteur identique $I$ de $\mathbf{C}$
dans $\mathbf{C}$, un foncteur résolvant de $I$ (appelé aussi \textit{résolution de
l'identité} dans $\mathbf{C}$) est donc un foncteur covariant \textit{exact} $C$ de $\mathbf{C}$
dans la catégorie des complexes à degrés positifs dans $\mathbf{C}$, muni d'une
augmentation $A\to C(A)$ qui soit un isomorphisme de $A$ sur
$Z^0C(A)$, et tel que $C(A)$ soit une résolution de $A$ si $A$ est injectif ;
alors $C(A)$ est une résolution de $A$ quel que soit $A$ en vertu de
proposition 2.5.1 (puisque le foncteur $I$ étant exact, on a $R^nI=0$ pour
$n>0$). Soient $C$ une résolution de l'identité dans $\mathbf{C}$, $F$ un
foncteur covariant exact à gauche de $\mathbf{C}$ cans $\mathbf{C}'$, \textit{supposons}
$R^nF(C^i(A))=0$ pour tout $n>0$ (i. e. les $C^i(A)$ $F$-acycliques).
\textit{Alors $\mathbf F(A)=F(C(A))$ est un foncteur résolvant pour $F$.} En
effet ce foncteur est exact, car si on a une suite exacte
$0\to A'\to A\to A''\to0$
on en conclut une suite exacte
$0\to C^i(A')\to C^i(A)\to
C^i(A'')\to0$ d'où, puisque $R^1F(C^i(A'))=0$, une suite exacte
$0\to F(C^i(A'))\to F(C^i(A))
\to F(C^i(A''))\to0$.
Ainsi (i) est vérifié ; il en est de même de (ii) puisque $F$ est exact à
gauche et $C$ est exact et de même de (iii) puisque on a ($C(A)$
étant une résolution $F$-acyclique de $A$)
$H^n(C(A))=R^nF(A)$, qui est nul si $n>0$ et si $A$ est injectif. --- Une façon
commode de construire une résolution de l'identité dans $\mathbf{C}$ est de prendre un
foncteur exact $C^0(A)$ de $\mathbf{C}$ dans $\mathbf{C}$ et un \textit{monomorphisme} fonctoriel
$A\to C^0(A)$ : on définit alors par récurrence les $C^n(A)$
($n\geqq0$) et les homomorphismes $d^{n-1}:C^{n-1}\to C^n$ en
prenant
$C^1(A)=C^0(C^0(A)/\operatorname{Im}A)$,
et, pour $n\geqq2$,
$C^n(A)=C^0(C^{n-1}(A)/\operatorname{Im}C^{n-2}(A))$,
$d^{n-1}$ étant défini à l'aide de l'homomorphisme d'augmentation de
$Q=C^{n-1}(A)/\operatorname{Im}C^{n-2}(A)$ dans $C^0(Q)$. On obtient bien
ainsi une résolution de l'identité ; pour que les $C^i(A)$ soient
$F$-acycliques, il suffit que les $C^0(A)$ le soient.

\LayerTarget{tohoku:example:2.5:hom-projective-resolution}
b) Soit $P=(P_n)$ une résolution projective d'un objet $A$ de $\mathbf{C}$.
Soit $F$ le foncteur $F(B)=\operatorname{Hom}(A,B)$ de $\mathbf{C}$ dans la catégorie
des groupes abéliens. Alors $F$ admet le foncteur résolvant
$\operatorname{Hom}(P,B)$.

\LayerTarget{tohoku:construction:resolving-functor-spectral-sequences:p1}
Nous allons maintenant montrer comment les suites spectrales les plus
importantes peuvent se calculer au moyen de foncteurs résolvants. Soit
$\mathbf F$ un foncteur résolvant d'un foncteur exact à gauche $F$ de $\mathbf{C}$
dans $\mathbf{C}'$. Soit $K$ un complexe dans $\mathbf{C}$, à degrés limités à gauche,
considérons le bicomplexe
$\mathbf F(K)=(\mathbf F^q(K^p))_{p,q}$, il dépend fonctoriellement de $K$, il
en est donc de même de ses deux suites spectrales et de leur aboutissement.
Les termes initiaux de ces deux suites spectrales s'écrivent aisément en
utilisant l'exactitude du foncteur $\mathbf F$ et la proposition 2.5.1, on
trouve
\LayerTarget{tohoku:eq:2.5.3}
\[
\tag{2.5.3}
I_2^{p,q}(\mathbf F(K))=H^p(R^qF(K))
\qquad
II_2^{p,q}(\mathbf F(K))=R^pF(H^q(K)).
\]
Ce sont là des isomorphismes fonctoriels avec les termes initiaux des suites

\TohokuContinuation{151}

